% !TeX root = ../../Book4.tex
% 纲要标题：### 7.2 映射锥
% 纲要标签：b4:sec:0602
% 纲要位置：计划纲要.md，第 3350 行。

\section{映射锥}
\label{b4:sec:0602}

一个复形映射也可以成为新复形的微分组成部分。映射锥将它在上同调上诱导的映射的核与余核信息集中到长正合列中，柱构造则把相关比较写成具体映射。

% 纲要内容（仅作写作计划，尚非教材正文）：
%
% * 锥与柱的显式构造
% * 锥的上同调长正合列
% * 移位和锥的符号检验
%

\subsection{锥与柱的显式构造}
\label{b4:subsec:0602:01}

上同调长正合列从短正合列出发，而许多问题最初只有一个映射 \(f:C\to D\)。要把它转成上同调之间的比较，可以构造一个逐次分裂的复形短列，左端为 \(D\)，右端为 \(C[1]\)。于是中项的次数 \(n\) 取 \(D^n\oplus C[1]^n=D^n\oplus C^{n+1}\)，两个对角微分分别取 \(\mathrm{d}_D^n\)、\(\mathrm{d}_{C[1]}^n=-\mathrm{d}_C^{n+1}\)。再把给定的 \(f^{n+1}:C^{n+1}\to D^{n+1}\) 放在非对角位置，便得到下面的映射锥。移位使这个非对角项的源、目标次数能够相接，所得长正合列将同时比较两端的上同调与 \(H(f)\)。

\begin{definition}\label{b4:def:mapping-cone}
设 \(f:C\to D\) 是加性范畴中的上链复形映射。规定
\[
\operatorname{Cone}(f)^n=D^n\oplus C^{n+1},\qquad
 \mathrm{d}_{\operatorname{Cone}(f)}^n=
 \begin{pmatrix}\mathrm{d}_D^n&f^{n+1}\\0&-\mathrm{d}_C^{n+1}\end{pmatrix}.
\]
所得复形称为 \(f\) 的\term[yingshe-zhui]{映射锥}[mapping cone]。在模中写作 \(\mathrm{d}(y,x)=(\mathrm{d}_Dy+f(x),-\mathrm{d}_Cx)\)。
\end{definition}
\symidx{\(\operatorname{Cone}(f)\)：复形映射的映射锥}

右下角的负号正是\cref{b4:def:complex-shift}中 \(C[1]\) 的微分符号。矩阵平方的非对角项为 \(\mathrm{d}_Df-f\mathrm{d}_C=0\)，对角项为两个微分的平方，因此确为复形。若把右下角的负号删去，非对角项通常会变成 \(\mathrm{d}_Df+f\mathrm{d}_C\)，不能保证为零。

映射柱要把 \(f\) 分解成 \(C\xrightarrow{j}\operatorname{Cyl}(f)\xrightarrow{r}D\)，其中 \(j\) 逐次分裂单，\(r\) 为同伦等价，并保持 \(rj=f\)。为给 \(C\) 留出嵌入，同时保留 \(D\) 的同伦信息，可以在 \(D\) 上添入由 \(C^n\oplus C^{n+1}\) 组成的可缩辅助复形。下面的三个分量正是 \(D\) 与这份辅助数据；定义后将用坐标变换识别它们。

\begin{definition}\label{b4:def:mapping-cylinder}
设 \(f:C\to D\) 是加性范畴中的上链复形映射。规定
\[
\operatorname{Cyl}(f)^n=D^n\oplus C^n\oplus C^{n+1},
\qquad
 \mathrm{d}(y,c,a)=(\mathrm{d}_Dy+f(a),\ \mathrm{d}_Cc-a,\ -\mathrm{d}_Ca).
\]
所得复形称为 \(f\) 的\term[yingshe-zhu]{映射柱}[mapping cylinder]。元组公式在一般加性范畴中表示相应双积矩阵。
\end{definition}
\symidx{\(\operatorname{Cyl}(f)\)：复形映射的映射柱}

这三项可以从一个坐标变换理解。令 \(z=y+f(c)\)，即把 \((y,c,a)\) 换成 \((z,c,a)\)，逆变换为 \((z,c,a)\mapsto(z-f(c),c,a)\)。因为 \(f\) 与微分交换，新坐标中的微分为
\[
\begin{gathered}
\mathrm{d}_Dy+f(a)+f(\mathrm{d}_Cc-a)
=\mathrm{d}_D(y+f(c))=\mathrm{d}_Dz,\\
\mathrm{d}(z,c,a)=(\mathrm{d}_Dz,\,\mathrm{d}_Cc-a,\,-\mathrm{d}_Ca).
\end{gathered}
\]
第一分量因而与后两分量分开：映射柱同构于 \(D\oplus K\)，其中 \(K^n=C^n\oplus C^{n+1}\)，微分为 \((c,a)\mapsto(\mathrm{d}_Cc-a,-\mathrm{d}_Ca)\)。辅助复形 \(K\) 可缩，收缩为 \((c,a)\mapsto(0,-c)\)。映射柱是在 \(D\) 上添入这一份同伦上为零的复形；同时 \(c\mapsto(0,c,0)\) 为 \(C\) 留出逐次分裂的嵌入。在新坐标中，该嵌入是 \(c\mapsto(f(c),c,0)\)，而到 \(D\) 的投影正是原坐标中的 \(y+f(c)\)。

\begin{proposition}\label{b4:prop:cylinder-factorization}
设 \(\mathcal A\) 为加性范畴，\(C,D\) 为其中的上链复形，\(f:C\to D\) 为复形映射。令 \(j:C\to\operatorname{Cyl}(f)\)、\(r:\operatorname{Cyl}(f)\to D\)、\(s:D\to\operatorname{Cyl}(f)\) 的逐次公式为
\[
j(c)=(0,c,0),\qquad r(y,c,a)=y+f(c),\qquad s(y)=(y,0,0).
\]
则 \(rj=f\)、\(rs=1_D\)，\(r,s\) 互为同伦逆。\(j:C\to\operatorname{Cyl}(f)\) 逐次为分裂单态射，且 \(q^n(y,c,a)=(y,a)\) 给出 \(j\) 的余核复形 \(\operatorname{Cone}(f)\)。在交换范畴中因而有逐次分裂短正合列
\[
0\longrightarrow C\xrightarrow{j}\operatorname{Cyl}(f)
\xrightarrow{q}\operatorname{Cone}(f)\longrightarrow0.
\]
\end{proposition}
\begin{proof}
将微分公式再用一次，三个分量分别为 \(\mathrm{d}_D^2y+\mathrm{d}_Df(a)-f\mathrm{d}_C(a)\)、\(\mathrm{d}_C^2c-\mathrm{d}_Ca+\mathrm{d}_Ca\)、\(\mathrm{d}_C^2a\)，均为零。由 \(f\) 与微分交换，\(j,r,s\) 都是复形映射；两个复合等式直接成立。

还需证明 \(1-sr\) 零伦。按同伦约定，比较恒等态射与 \(sr\)，要找到 \(h\) 使 \(1-sr=\mathrm{d}h+h\mathrm{d}\)。定义降次态射 \(h(y,c,a)=(0,0,-c)\)，则
\[
\mathrm{d}h(y,c,a)=(-f(c),c,\mathrm{d}_Cc),\qquad
h\mathrm{d}(y,c,a)=(0,0,-\mathrm{d}_Cc+a),
\]
相加得 \((-f(c),c,a)=(1-sr)(y,c,a)\)。因此 \(sr\simeq1\)。

\(j^n\) 的逐次左逆是 \(\pi_C^n(y,c,a)=c\)，满足 \(\pi_C^nj^n=1_{C^n}\)。逐次投影到第一、第三个双积分量得到 \(q\)，其逐次截面为 \(\sigma^n(y,a)=(y,0,a)\)，满足 \(q^n\sigma^n=1\)。其目标上的微分为 \((\mathrm{d}_Dy+f(a),-\mathrm{d}_Ca)\)，正是映射锥的微分，且 \(qj=0\)。若复形映射 \(u:\operatorname{Cyl}(f)\to E\) 满足 \(uj=0\)，则双积泛性质给出唯一的态射族 \(v^n:\operatorname{Cone}(f)^n\to E^n\)，使 \(u^n=v^nq^n\)。由于每个 \(q^n\) 都是分裂满态射，等式
\[
\begin{aligned}
\bigl(\mathrm{d}_E^nv^n-v^{n+1}\mathrm{d}_{\operatorname{Cone}(f)}^n\bigr)q^n
&=\mathrm{d}_E^nu^n-u^{n+1}\mathrm{d}_{\operatorname{Cyl}(f)}^n\\
&=0
\end{aligned}
\]
可以在右边消去满态射 \(q^n\)，得到 \(\mathrm{d}_E^nv^n=v^{n+1}\mathrm{d}_{\operatorname{Cone}(f)}^n\)。所以逐次分解组成复形映射 \(v\)，\(q\) 满足 \(j\) 在复形范畴中的余核泛性质。
\end{proof}

这里逐次分裂是指每个次数上的短正合列分裂；复形分裂还要求分裂态射与微分交换。上面的 \(\pi_C\) 与 \(\sigma\) 通常不满足这个要求：\(\pi_C\mathrm{d}(y,c,a)=\mathrm{d}_Cc-a\)，而 \(\mathrm{d}_C\pi_C(y,c,a)=\mathrm{d}_Cc\)；\(\mathrm{d}\sigma(y,a)\) 的中间分量为 \(-a\)，而 \(\sigma\mathrm{d}(y,a)\) 的中间分量为零。可缩则是单个复形的恒等映射零伦的条件；这里可缩的是辅助复形 \(K\)，映射柱本身仍保留 \(D\) 的同伦信息。

\subsection{锥的上同调长正合列}
\label{b4:subsec:0602:02}

\begin{proposition}\label{b4:prop:cone-long-exact}
设 \(\mathcal A\) 是交换范畴，\(C,D\) 为其中的上链复形，\(f:C\to D\) 是复形映射。嵌入 \(y\mapsto(y,0)\) 与投影 \((y,x)\mapsto x\) 给出逐次分裂短正合列
\[
0\longrightarrow D\longrightarrow\operatorname{Cone}(f)
\longrightarrow C[1]\longrightarrow0.
\]
其连接态射在 \(H^n(C[1])=H^{n+1}(C)\) 的识别下就是 \(H^{n+1}(f)\)。因此有自然上同调长正合列
\[
\cdots\to H^n(C)\xrightarrow{H^n(f)}H^n(D)
\to H^n\operatorname{Cone}(f)\to H^{n+1}(C)
\xrightarrow{H^{n+1}(f)}H^{n+1}(D)\to\cdots.
\]
\end{proposition}
\begin{proof}
微分矩阵表明嵌入与投影均为复形映射，前者为后者的核，后者为前者的余核；这些泛性质由逐次双积分解得到。先应用\cref{b4:thm:cohomology-long-exact}，所得原始片段为
\[
H^n(D)\longrightarrow H^n\operatorname{Cone}(f)
\longrightarrow H^n(C[1])\xrightarrow{\delta^n}H^{n+1}(D).
\]
再用 \(H^n(C[1])=H^{n+1}(C)\) 识别移位项，并计算 \(\delta^n\)。投影的逐次截面 \(x\mapsto(0,x)\) 不必为复形映射，因为其微分还产生 \(f(x)\) 分量。对余圈子对象 \(Z^{n+1}(C)\) 使用这一截面，再作锥微分，得到进入 \(D^{n+1}\) 的态射 \(f^{n+1}|_{Z^{n+1}(C)}\)。由\cref{b4:prop:connecting-map-construction}，它在上同调上诱导的正是连接态射 \(H^{n+1}(f)\)。在模中，取余圈 \(x\in Z^{n+1}(C)\)，它代表 \(H^n(C[1])\) 中的同一个类。计算为
\[
\begin{aligned}
x&\longmapsto(0,x),\\
\mathrm{d}_{\operatorname{Cone}(f)}^n(0,x)
&=(f^{n+1}(x),-\mathrm{d}_C^{n+1}x)\\
&=(f^{n+1}(x),0),
\end{aligned}
\]
故 \(\delta^n[x]=[f^{n+1}(x)]\)，没有额外负号。最终长列中的前一个箭头 \(H^n(C)\to H^n(D)\) 来自上一次数的连接 \(\delta^{n-1}:H^{n-1}(C[1])\to H^n(D)\)，在移位识别下正是 \(H^n(f)\)。
\end{proof}

\begin{example}\label{b4:exm:cone-single-map}
把左 \(R\)-模 \(M,N\) 放在次数零，同态 \(u:M\to N\) 的锥只有次数 \(-1,0\) 非零，恰为 \(M\xrightarrow{u}N\)。所以
\[
H^{-1}\operatorname{Cone}(u)=\ker u,\qquad
H^0\operatorname{Cone}(u)=\operatorname{coker}u.
\]
例如 \(\mathbb Z\xrightarrow{m}\mathbb Z\) 的锥在 \(m\ne0\) 时只剩零次上同调 \(\mathbb Z/m\mathbb Z\)。因此 \(\operatorname{Cone}(u)\) 零调，当且仅当 \(\ker u=0\)、\(\operatorname{coker}u=0\)，也就是 \(u\) 为模同构。下一节的\cref{b4:prop:cone-quasi-isomorphism}将把这一原型推广到任意复形映射：锥零调恰好检测它是否在所有次数诱导上同调同构。
\end{example}

\subsection{移位和锥的符号检验}
\label{b4:subsec:0602:03}

锥的不同符号约定可以彼此同构，但必须明确写出比较映射。例如有些文献把 \(C^{n+1}\) 放在前面，并在 \(f\) 前加负号；这时交换两个分量并改变其中一个分量的符号，才能与本书定义比较。

\begin{proposition}\label{b4:prop:cone-sign-compatibility}
设 \(\mathcal A\) 为加性范畴，\(C,D\) 为其中的上链复形，\(f,g:C\to D\) 为复形映射。以下元组公式表示相应双积上的态射矩阵。
\begin{enumerate}
\item 若态射族 \(h^n:C^n\to D^{n-1}\) 满足 \(f-g=\mathrm{d}_Dh+h\mathrm{d}_C\)，则次数 \(n\) 上的公式 \((y,x)\mapsto(y+h^{n+1}(x),x)\) 给出 \(\operatorname{Cone}(f)\cong\operatorname{Cone}(g)\)。这个同构由所给的 \(h\) 确定，不同同伦一般会给出不同同构。
\item 对整数 \(r\)，映射 \((y,x)\mapsto(y,(-1)^r x)\) 给出 \(\operatorname{Cone}(f[r])\cong\operatorname{Cone}(f)[r]\)。
\item \(\operatorname{Cone}(1_C)\) 可缩，收缩同伦为 \((y,x)\mapsto(0,y)\)。
\end{enumerate}
\end{proposition}
\begin{proof}
第一项的映射与微分交换等价于
\[
\mathrm{d}_Dy+\mathrm{d}_Dh(x)+g(x)=\mathrm{d}_Dy+f(x)-h(\mathrm{d}_Cx),
\]
这恰为同伦公式；逆映射把 \(h\) 换成 \(-h\)。第二项中 \(f[r]^n=f^{n+r}\)，不另加符号。源 \(\operatorname{Cone}(f[r])^n\) 与目标 \(\operatorname{Cone}(f)[r]^n\) 的底对象都是 \(D^{n+r}\oplus C^{n+r+1}\)，差别只在微分。暂略已经移到 \(n+r\) 的指标，源的微分矩阵为
\[
\begin{pmatrix}(-1)^r\cdot\mathrm{d}_D&f\\0&-(-1)^r\cdot\mathrm{d}_C\end{pmatrix},
\]
目标的微分矩阵为 \((-1)^r\begin{pmatrix}\mathrm{d}_D&f\\0&-\mathrm{d}_C\end{pmatrix}\)。令 \(\epsilon=(-1)^r\)，比较态射 \(t^n=\operatorname{diag}(1,\epsilon)\) 在第二分量补上符号。用 \(\epsilon^2=1\) 计算，两种复合的矩阵均为
\[
t^{n+1}\mathrm{d}_{\operatorname{Cone}(f[r])}^n
=\mathrm{d}_{\operatorname{Cone}(f)[r]}^nt^n
=\begin{pmatrix}\epsilon\cdot\mathrm{d}_D&f\\0&-\mathrm{d}_C\end{pmatrix}.
\]
而 \((t^n)^2=1\)，故给出复形同构。第三项令 \(h(y,x)=(0,y)\)，则 \(\mathrm{d}h(y,x)=(y,-\mathrm{d}y)\)、\(h\mathrm{d}(y,x)=(0,\mathrm{d}y+x)\)，其和为 \((y,x)\)。
\end{proof}

移位后的同伦也须改符号：若 \(h\) 是 \(f\simeq g\) 的同伦，则 \(h_r^n=(-1)^r h^{n+r}\) 是 \(f[r]\simeq g[r]\) 的同伦。这从两边微分均乘 \((-1)^r\) 直接得出。当 \(r\) 为奇数时，若只移动指标而不乘这个符号，同伦公式右边会多出整体负号；\(r\) 为偶数时则无此差异。
